\uuid{zPPh}
\exo7id{3810}
\titre{exo7 3810}
\auteur{quercia}
\organisation{exo7}
\datecreate{2010-03-11}
\isIndication{false}
\isCorrection{true}
\chapitre{Espace euclidien, espace normé}
\sousChapitre{Problèmes matriciels}
\cpge{mpsi-mp2i.calcul-matriciel-systemes.e, pcsi.determinants.c, ptsi.determinants.a, mpsi-mp2i.changements-bases.c, mp-mpi.reduction-endomorphismes.c, pc.reduction.b, psi.reduction-endomorphismes.b, pt.reduction-endomorphismes.b, mp-mpi.endomorphismes-euclidien.b, pc.endomorphismes-euclidiens.b, psi.endomorphismes-euclidien.b, pt.isometries-euclidien.b, pc.complements-espaces-vectoriels.c, psi.complements-espaces-vectoriels.c, pt.complements-algebre-lineaire.c, bcpst2.valeurs-propres-vecteurs-propres.a, bcpst2.valeurs-propres-vecteurs-propres.b}
\module{Algèbre}
\niveau{L2}
\difficulte{}

\contenu{
\question{
Soit $A \in \mathcal{M}_n(\C)$ de valeurs propres $\lambda_1, \dots, \lambda_n$.
Montrer que $AA^* = A^*A \iff \mathrm{tr}(AA^*) = |\lambda_1|^2+ \dots + |\lambda_n|^2$.
}
\reponse{
trigonaliser $A$ dans une base orthonormée.
}
}
